% SGA 3, Expose VII_B, English reconstruction -- Loop 1.
% Scope: Section 5.2.5 in full on authority local pages 94--95;
% stops immediately before Section 5.2.6.
% Combined-reader pages: 607--608. Re-edition printed pages: 596--597.
% Original-edition marginal locator in scope: 562.
% Editorial notes in scope: 162--163.
% Authority: Polo--Gille Exp7B-13oct24.pdf,
% SHA-256 D303690E0195654BBE7C3DAA5D89EA23E24CD709A3575217234AC19AE8836694.
% Jacob Reinhold's CC BY 4.0 Markdown was consulted as comparison lineage.
% No genuine diagram occurs in this component; all mathematics is native LaTeX.

\subsubsection*{5.2.5}
\phantomsection
\label{sga3:VIIB:5.2.5}

还需考虑 \(B\) 高度无限，且投影 \(I\to\omega\) 有与 \(\omega\) 和 \(I\) 的滤过相容的截面 \(\sigma\) 的情形。考虑由 \(\sigma\) 诱导的态射：
\[
  \phi:
  B'
  =
  k[[\omega]]
  \big/
  \bigl((x^{p^n})\bigr)_{x\in\omega_n}
  \longrightarrow B.
\]
只须证明，对每个 \(r\in\overline{\mathbb N}\)，由 \(\phi\) 诱导的映射
\[
  \phi_r:B'_r\longrightarrow B_r
\]
为同构。\footnote{%
\label{sga3:VIIB:nde:162}事实上，\(B'\) 和 \(B\) 分别为各 \(B'_r\) 和各 \(B_r\) 的逆极限。此外，为顾及 \SGARefLink{sga3:VIIB:5.2.3}{5.2.3} 中所作的增补，编者修改了原文以下部分。}

对每个 \(r\in\mathbb N^*\)，令 \(G_r\) 为迭代 Frobenius 态射
\[
  G\longrightarrow G^{(p^r)},
\]
的核，令 \(HG_r\) 为 \(H^{(p^r)}\subset G^{(p^r)}\) 在 \(G\) 中的逆像形式子群。因此 \(HG_r\) 由 \(AI\) 的元素的 \(p^r\) 次幂生成的闭理想定义，即
\[
  A\check J_r,
  \qquad
  \check J_r=\{\,x^{p^r}\mid x\in I\,\}.
\]
完全如同 \SGARefLink{sga3:VIIB:5.2.3}{5.2.3}，可知
\[
  B_r=B/B\check J_r
\]
为 \(HG_r/H\) 的仿射代数；特别地，其高度至多为 \(r\)。

令
\[
  \mathfrak m(r)=I/B\check J_r
\]
为 \(B_r\) 的增广理想。典范投影 \(B\to B_r\) 诱导同构
\[
  \omega=I/I^2
  \xrightarrow{\ \sim\ }
  \omega(r)
  =
  \mathfrak m(r)/\mathfrak m(r)^2,
\]
我们用此同构将这两个空间等同。令 \(\omega(r)_n\) 为闭理想
\[
  \mathfrak m(r)_n
  =
  \{\,y\in\mathfrak m(r)\mid y^{p^n}=0\,\}.
\]
在 \(\omega(r)\) 中的像。它也是闭理想
\[
  I(r)_n
  =
  \{\,x\in I\mid x^{p^n}\in B\check J_r\,\}.
\]
在 \(\omega\) 中的像。
\footnote{\label{sga3:VIIB:nde:163}编者扩充了原文以下部分。}

显然，
\[
  \omega(r)_n=\omega
  \qquad\text{若 }n\geq r.
\]
我们将证明
\[
  \omega(r)_n=\omega_n
  \qquad\text{若 }n<r.
\]
对每个 \(r,n\)，序列
\[
  0
  \longrightarrow
  I(r)_n\cap I^2
  \longrightarrow
  I(r)_n
  \longrightarrow
  \omega(r)_n
  \longrightarrow0
\]
正合。对固定的 \(n\)，有
\[
  \bigcap_r I(r)_n=I_n,
\]
因为 \(\bigcap_rB\check J_r=0\)。在 \(\operatorname{PC}(k)\) 中，滤过逆极限是正合的\textup{（参见 \SGARefLink{sga3:VIIB:0.2}{0.2}）}，所以对每个 \(n\)，
\[
\tag{*}\label{sga3:vii:remark:5-2-4-a:equation:star}
  \omega_n=\bigcap_r\omega(r)_n.
\]
另一方面，\SGARefLink{sga3:VIIB:5.2.4.A}{注~5.2.4.A} 给出
\[
  \omega(r)_n=\omega(r+1)_n
  \qquad\text{若 }n<r.
\]
结合 \textup{(*)}，即得
\[
  \omega(r)_n=\omega_n
  \qquad\text{若 }n<r.
\]

因此，按子空间 \(\bigl(\omega(r)_n\bigr)_{n\geq0}\) 赋予滤过的向量空间 \(\omega\)，恰好是 \SGARefLink{sga3:VIIB:5.2:theorem}{5.2} 的记号中的 \({}^{r}\omega\)。所以复合
\[
  \sigma(r):
  \omega\xrightarrow{\ \sigma\ }I
  \longrightarrow\mathfrak m(r)
\]
与滤过 \(\bigl(\omega(r)_n\bigr)\) 和 \(\bigl(\mathfrak m(r)_n\bigr)\) 相容。由于
\[
  k[[{}^{r}\omega]]
  \big/
  \bigl((x^{p^n})\bigr)_{x\in{}^{r}\omega_n}
\]
恰为 \(B'_r\)，且 \(\phi_r:B'_r\to B_r\) 由 \(\sigma(r)\) 诱导，已经对有限高度代数证明的结果表明 \(\phi_r\) 为同构。
